%% kapa: exam=1 round=33 year=2022 period=2 form=A subject=appraisal_law subject_ko=감정평가관계법규 seq=2 number=2
%% source: Q-Net 감정평가사 기출문제 (cst00309) artlSeq=5209997 file=raw/1008/5209997/2022년도 제33회 감정평가사 1차시험 2교시(감정평가 관계법규, 회계학).pdf
%% conversion: kordoc 4.24.51 -> tools/gen1.py (auto); review=proofread-sonnet; flags=none
\begin{question}{2}
국토의 계획 및 이용에 관한 법령상 개발행위에 따른 공공시설 등의 귀속에 관한 설명으로 옳지 않은 것은?
\begin{choices}
  \item 개발행위허가를 받은 자가 행정청인 경우 개발행위허가를 받은 자가 새로 공공시설을 설치한 경우 새로 설치된 공공시설은 그 시설을 관리할 관리청에 무상으로 귀속된다.
  \item 개발행위허가를 받은 자가 행정청인 경우 개발행위허가를 받은 자가 기존의 공공시설에 대체되는 공공시설을 설치한 경우 종래의 공공시설은 개발행위허가를 받은 자에게 무상으로 귀속된다.
  \item 개발행위허가를 받은 자가 행정청이 아닌 경우 개발행위허가를 받은 자가 새로 설치한 공공시설은 그 시설을 관리할 관리청에 무상으로 귀속된다.
  \item 개발행위허가를 받은 자가 행정청이 아닌 경우 개발행위로 용도가 폐지되는 공공시설은 개발행위허가를 받은 자에게 무상으로 귀속된다.
  \item 특별시장·광역시장·특별자치시장·특별자치도지사·시장 또는 군수는 공공시설의 귀속에 관한 사항이 포함된 개발행위허가를 하려면 미리 관리청의 의견을 들어야 한다.
\end{choices}
\end{question}
\answer{4}
